Aus $F = G\,\frac{m_1 m_2}{r^2}$ folgt für das Verhältnis der Kräfte bei den Abständen $r_1$ und $r_2$
\begin{align}
 \frac{F_2}{F_1} &= \left(\frac{r_1}{r_2}\right)^2 = k
\end{align}
mit $k = \kfO$. Aufgelöst nach dem Verhältnis der Abstände:
\begin{align}
 \frac{r_2}{r_1} &= \rrF = \frac{1}{\sqrt{\kfO}} = \rr \approx \result{\rrP{0}{2}}
\end{align}
Der Abstand muss also um den Faktor $\sqrt{3}$ verkleinert werden, d.\,h. um
\begin{align}
 \frac{r_2-r_1}{r_1} &= \rr-1 \approx \result{\drP{0}{2}}
\end{align}
